%2025-03-22, from the discussion with DEEPSEEK (mainly) and JULIUS AI, knowing that Lemma 8 in "A Proof of the Riemann %Hypothesis Based on a New Expression of the Completed Zeta Function" can be extended to some kinds of L-functions, %i.e., the GRH can be proved with some modifications.

%2025-06-01 modified from "Extending result..."
%2025-07-27/2025-07-30
%2025-09-10 修改公式24：漏掉了一项(-1)的N次方，这个遗漏不影响定理4的结果。
%2025-09月, V7主要是考虑一致收敛性
%2025-10月国庆假期，V8主要修改定理1以及引理8的证明（适配引理5的条件，之前引理8的证明啰嗦且有错）
%2025-10-09 two modifications before and in the proof of Theorem 1 to emphasize the zeros in quadruple. 还是版本V8
%2025-10-11 The critical strip has been extended from $(0, 1)$ to $[0, k]$, which is of very important significance; V9 %2025-10-11 Theorem 10 and Lemma 10 are added. V9
%2025-10-14 梳理了定理3-4的条件，特别是在定理5的证明中，明确了1）-5）的条件，方便后面定理5-7的证明中逐一对应。V10
%2025-12-05 去掉一致收敛性条件。V11
%2026-03-12,03-18 不依赖PREPRINT论文的独立发表版本。去掉了之前的引言/附录，增加了“Preliminary Lemmas"，在DEDEKIND ZETA
%FUNCTION的定义后面，以及Dirichlet$L$-function的定义后面加了黎曼zeta函数是其特例的说明，润色了摘要，修改了几处TYPO, V13
%2026-03-19 今天晚上发现V9将critical strip extended from $(0, 1)$ to $[0, k]$是不必要的，且在 Theorem4    %的证明中以及后来MAIN RESULT中的相关内容带来了不必要的麻烦和啰嗦，所以去掉这一扩展，涉及THEOREM 1-10
%2026-03-21 今天上午发现2026-03-19的修改不对，因为我们的依据是LEMMA 5.5:其给出的临界带是$[0, k]$,V14
%2026-03-30 今天修改论文题目 V14 Details 0f changes: 1. The title of the paper is modified to be more specific; 2. The %2026-03-30 subsection Modular Form L-Function is removed; 3. Some refinements in the proofs of Theorem 3-Theorem 8
%2026-03-31 One could refer to the web site of The American Institute of Mathematics for the statements of ERH, GRH, %2026-03-31 and Grand Riemann Hypothesis: https://www.aimath.org/WWN/rh/,
%2026-03-31 https://www.aimath.org/WWN/rh/articles/html/18a/, https://www.aimath.org/WWN/rh/articles/html/17a/, %2026-03-31 https://www.aimath.org/WWN/rh/articles/html/19a/
%2026-04-02 在定理5的证明前面加了一个REMARK，填上了一个逻辑上的GAP，定理5和定理7的证明中CASE 1的说明进行了简化，V15
%2026-04-03 在L(\chi,s)的完备化中，限定\chi is primitive；继续简化THEOREM 5证明之前的REMARK，3个合并成两个. V15
%2026-04-16 两处修改，定理4的证明之前，加了一个REMARK；修改了grand RH的版本1，消除逻辑gap，V15
%2026-04-25 今天发现LEMMA 5.5的闭区间临界带，可以结合SARNAK的结果，修正为开区间，这样可以做很多简化，涉及67+4(Theorem %2026-04-25 1-4)+8(Theorem 5-9),总共约80处修改。
%2026-04-26 Theorem 1-Theorem 2 的收敛条件修改，原理是关于\rho_i的，改为关于\beta_i
%2026-04-27 %最近反复确认：整除性相关的结果并不需要限制零点的位置，也就是不需要限制零点在临界带之内(包括边界)，需要限制零点在临界带以内的是保证整除性存在的条件$\sum_{i=1}^{\infty}\frac{1}{\beta_i^2}<\infty$
%2026-05-07 Theorem 1 和 Theorem 2的收敛条件进行了修正 V16
%2026-06-03 于波士顿蔡兴敏家，继续修改V16。the proof of Theorem 1 has been changed.
%2016-06-15 修改DEDEKIND ZETA FUNCTION有关内容，主要是定义部分，COMPLETED部分内容；自守形式L函数有关内容，主要是FUNCTIONAL EQUATION部分，以及后面THEOREM 7的证明部分。1.The proof of Theorem 1 has been modified, Eq.(7b) and Eq.(7c) have been added; 2.Eq.(30) has been modified; 3.The Definition of the Dedekind zeta function has been modified, the completed Dedekind zeta function Eq. (38) has been modified, the functional equation Eq.(39) has been modified; 4. Eq.(45), Eq.(47), and Eq.(48) have been modified.
%2026-07-07 from absolute convergence to absolute and uniform convergence; the proof of Theorem 6.
%1. The paper title is modified; 2. The Abstract is modified; 3. The proof of Lemma 2 is modified; 4. The proof of Lemma 3 is modified; 5. The "absolute convergence" has been changed to "absolute and uniform convergence" throughout the paper, affected Lemma 2, Theorem 1, Theorem 2; 6. The proof of Theorem 6 is modified; 7. A new reference is added.
%2026-09-11 in the proof of Theorem 7/CASE 2, added two points(two sentences) according to the proof of Theorem 5; and in the proof of Theorem 5, added a paragraph according to the proof of Theorem 7.
% V20：修改了题目，摘要，关键词，修改了SKETCH OF THE PROOF;修改了定理1证明中对称方程导出整除性的论述部分，增加了MAIN
%RESULTS里面的GENERALPROPERTY论述部分；修改了公式（30）及W(\chi)的意义；修改了公式（29）下面的REMARK中的q=1;三种L函数定义前面都加了一句说明，
%并且调整了AUTOMORPHIC L-函数的定义：把自守形式的说明放在了定义中；修改了致谢部分。2026-09-30
%2026-10-04，修改了定理2的表达式EQ.(8)；调整了参考文献及顺序，去掉了自己的PREPRINT论文及引用；在MAIN %RESULTS定理5证明之前增加了一段说明，解释了主特征情况的处理办法;修改了定理1-4的表述；
%2026-10-09：1.Modified Eq. (8) in Theorem 2.Adjusted the references and their order, and removed one of %reference.3.Added a
% paragraph before the proof of Theorem 5 in MAIN RESULTS, explaining the treatment of the principal character case.
% Revised the statements of Theorems 1–4.
% 2026-10-17: TBD: 1, 改进THEOREMS 1-4的表达-finished； 2. THEOREMS 5-6补上欧拉乘积-finished,10月21日; m_i in Lemma 3 finished;
% m_i的性质全文都没有论述-论文最前面有一句交代，应该补充两个引理; 是否应该强调polynomials with undetermined %coefficients？不用。
% 2026-10-28: Lemma 4 and Lemma 5 are added; Theorem 3 and Theorem 4 are modified with \overline\lambda replaced;
% 2026-1--29: Theorems 1-4 are modified with statements;
% 2026-10-31: Two remarks are added, one after Theorem 1, one before Theorem 3.
\documentclass[12pt]{article}
\usepackage{amsmath, amsthm, amssymb, amsfonts}
\usepackage[utf8]{inputenc}
\usepackage{natbib}
\usepackage{array}
\usepackage{graphicx}


\title{On the Extended, Generalized, and Grand Riemann Hypotheses: A Unified Approach via the Divisibility Property of Entire Functions}
\author{Weicun Zhang \thanks{Corresponding author: weicunzhang@ustb.edu.cn}\\ University of Science and Technology Beijing\\Beijing 100083, China\\\\
* Previous versions of this manuscript are available at \\https://www.preprints.org/manuscript/202506.0481}
\date{\today}
\newtheorem{theorem}{Theorem}
\newtheorem{lemma}{Lemma}
\newtheorem{definition}{Definition}
\renewcommand{\proofname}{Proof}

\begin{document}

\maketitle

\begin{abstract}

 This paper addresses the Extended, Generalized, and Grand Riemann Hypotheses within a unified framework via the divisibility property of entire functions (see Lemma 2), based on general properties of \(L\)-functions (see P1-P6). First, pairing complex-conjugate zeros in the Hadamard products of the completed \(L\)-functions \(\Lambda(\lambda, s)\) (for the self-dual case) or \(\Lambda(\lambda, s)\Lambda(\bar\lambda, s)\) (for the non-self-dual case), together with the fact that all zeros lie in the critical strip, as guaranteed by the Euler product property, yields absolutely and uniformly convergent infinite products with irreducible real quadratic polynomial factors. Second, substituting these infinite products into the symmetric functional equation yields the divisibility property of the completed \(L\)-functions. Third, this divisibility property, combined with the uniqueness of zero multiplicities, forces all zeros of \(\Lambda(\lambda, s)\) or \(\Lambda(\lambda, s)\Lambda(\bar\lambda, s)\) in the critical strip to lie on the critical line. Consequently, the existence of Landau–Siegel zeros is excluded, thereby confirming the Landau–Siegel zeros conjecture.

As for the Davenport--Heilbronn counterexample, since it possesses no Euler product---the fundamental structural property that confines zeros to the critical strip, thereby guaranteeing the absolute and uniform convergence of the infinite products and, in turn, the divisibility property of entire functions---it falls outside the scope of the methods and conclusions of this paper.\\

\textbf{MSC code:} 11M26\\

\textbf{Keywords:} Extended Riemann Hypothesis, Generalized Riemann Hypothesis, Grand Riemann Hypothesis, Landau–Siegel zeros conjecture, Davenport--Heilbronn counterexample, divisibility property of entire functions, uniqueness of zero multiplicities, symmetric functional equation, Hadamard product, Euler product.
\end{abstract}

\section{Notation and Sketch of the Contents}

\textbf{Notation.}\\

\noindent $\mathbb{N}$: the set of natural numbers, i.e., $\mathbb{N}=\{1, 2, 3, \ldots\}$.\\
$\mathbb{R}$: the field of real numbers.\\
$\mathbb{C}$: the field of complex numbers (the complex plane).\\
$\mathbb{Q}$, the field of rational numbers.\\
$\mathbb{R}[x]$: the ring of real polynomials.

It should be noted that in this paper, '$j$' is used to denote the imaginary unit ($j^2=-1$), while '$i$' together with '$l, n, t$' serves as natural number indices.

It is also worth noting that the following fact (supported by Lemma 4 and Lemma 5) plays an important role in this paper: \textbf{for a given non-zero entire function, the multiplicities of its zeros are finite and uniquely determined, hence invariant, even though their specific values may be unknown.}\\

\textbf{Sketch of the Contents.}\\

\noindent\textbf{Preliminary Lemmas}: Five Lemmas are provided to support the content that follows.\\
\textbf{Foundational Theorems}: Based on Lemma 2, four theorems are progressively proved from Theorem 1 to Theorem 4.\\
\textbf{Main Results}: Based on Theorem 4, the generalized Riemann Hypothesis (for Dirichlet $L$-functions) is proved in Theorem 5; the extended Riemann Hypothesis (for Dedekind zeta functions) is proved in Theorem 6; Theorem 7 deals with the non-trivial zero distribution of cuspidal automorphic $L$-functions, which corresponds to grand Riemann Hypothesis (version 1); Theorem 8 deals with the non-trivial zero distribution of all kinds of $L$-functions, which corresponds to grand Riemann Hypothesis (version 2).


\section{Preliminary Lemmas}

\textbf{Lemma 1:} Let  $m(x), g_1(x), . . ., g_n(x) \in \mathbb{R}[x], n\geq 2$. If $m(x)$ is irreducible (prime) and divides the product $g_1(x) \cdots g_n(x)$, then $m(x)$ divides one of the polynomials $g_1(x), . . . , g_n(x)$.

\textbf{Remark:} The contents of Lemma 1 can be found in many textbooks of linear algebra, modern algebra,  or abstract algebra, such as Ref.[1, p.135].

Next, we extends the above polynomial divisibility property in Lemma 1 to entire functions expressed as absolutely and uniformly convergent infinite products with irreducible real polynomial factors.

\textbf{Lemma 2:} Let $f(s)$ be an entire function on $\mathbb{C}$ with $f(0) \neq 0$.
Suppose $f(s)=\prod_{i=1}^{\infty} g_i(s)$ converges absolutely and
uniformly on compact subsets of $\mathbb{C}$, where each
$g_i(s)\in\mathbb{R}[s]$ is irreducible of degree
$d\in\{1,2\}$. If $m(s)\in\mathbb{R}[s]$ is irreducible
of degree $d$ and divides $f(s)$, then $m(s)$ divides $g_i(s)$ for some
$i\geq 1$.


\textbf{Proof:}
Let \(\alpha\) be a root of \( m(s) \), i.e., \( m(\alpha) =0 \). Since \( m(s) \mid f(s) \), we have \( f(\alpha) = 0 \). By absolute convergence of \(\prod_{i=1}^{\infty} g_i(s)\) at $s=\alpha$, there exists at least one index \( i \in \mathbb{N} \) such that \( g_i(\alpha) = 0 \), otherwise \( g_i(\alpha) \neq 0 \) for all $i$, then \(\prod_{i=1}^{\infty} g_i(\alpha)\) converges to a non-zero limit, contradicting \( f(\alpha) = 0 \).

As \( g_i(s) \) and \( m(s) \) are irreducible in \(\mathbb{R}[s]\) with \(\deg(g_i(s)) = \deg(m(s)) = d\), they share the root \(\alpha\). Thus:
\begin{itemize}
  \item If \( d = 1 \), then \( g_i(s) = a(s - \alpha) \) and \( m(s) = b(s - \alpha) \) for \( a, b \neq 0 \), so \( m(s) \mid g_i(s) \).
  \item If \( d = 2 \), then both \( g_i(s) \) and \( m(s) \) have roots \(\{\alpha, \bar\alpha\}\), so \( g_i(s) = c \cdot m(s) \) for \( c \neq 0 \), hence \( m(s) \mid g_i(s) \).
\end{itemize}
In both cases, \( m(s) \) divides \( g_i(s) \).

That completes the proof of Lemma 2.\\

\textbf{Remark:} Although the uniform convergence is not used directly in the proof of Lemma 2, it guarantees that the infinite product defines an entire function.


\textbf{Lemma 3:} The infinite product $\prod_{i=1}^{\infty} \Big{(}\frac{\beta_i^2}{\alpha_i^2 + \beta_i^2}\Big{)}^{m_i}$
converges to a non-zero constant, where $ 0\leq \alpha_i \leq k, 0<|\beta_1| \leq|\beta_2|\leq|\beta_3|, \cdots, \sum_{i=1}^{\infty} \frac{1}{\beta_i^2} < \infty$ (with each
$\beta_i^2$ appearing $m_i$ times), and $1\leq m_i< \infty$ is the multiplicity of zero $\rho_i=\alpha_i + j\beta_i$ of $\xi(s)$.

\textbf{Proof:} Let $a_i = \frac{\alpha_i^2}{\alpha_i^2+\beta_i^2}$, then
$\frac{\beta_i^2}{\alpha_i^2+\beta_i^2} = 1-a_i$.  Since \( 0 \leq \alpha_i \leq k \) and \( \beta_i \neq 0 \), we have $0 \leq a_i \leq \frac{k^2}{\beta_i^2}$.
By the given condition $\sum_{i=1}^{\infty}\frac{1}{\beta_i^2}<\infty$ (with each
$\beta_i^2$ appearing $m_i$ times), we have equivalently
$\sum_{i=1}^{\infty}\frac{m_i}{\beta_i^2}<\infty$ (with $|\beta_1|<|\beta_2|<\cdots$).
Hence $\sum_{i=1}^{\infty}m_i a_i \leq k^2\sum_{i=1}^{\infty}\frac{m_i}{\beta_i^2}< \infty$, which guarantees that the product
$\prod_{i=1}^{\infty}(1-a_i)^{m_i} = \prod_{i=1}^{\infty}\left(\frac{\beta_i^2}{\alpha_i^2+\beta_i^2}\right)^{m_i}$
converges to a non-zero constant.

That completes the proof of Lemma 3.\\

\textbf{Lemma 4:} Let $f(s)$ be a non-zero entire function, and let $s_0$ be a zero of $f(s)$. Then the multiplicity of $s_0$ is a finite positive integer.

\textbf{Proof:} Let $f(s) \not\equiv 0,  s \in \mathbb{C}$,  be an entire function, which means it is holomorphic  on the whole complex plane. Suppose $f(s)$ has a zero at $s_0 \in \mathbb{C}$ of multiplicity $m$, then $ f(s) = (s - s_0)^m g(s) $, where $g(s)$ is also an entire function and $g(s_0) \neq 0$.

Assume for contradiction that $m$ is infinite, which implies there exists an accumulation point of zeros in the neighbor of $s_0$.  Then, by Identity Theorem for holomorphic functions, and considering "0" is also an entire function, we have $f(s) \equiv 0, s \in \mathbb{C}$, which contradicts the given condition that $f(s)\not\equiv 0, s \in \mathbb{C}$.
Thus, the assumption is false, i.e., $m$ must be a finite positive integer.

That completes the proof of Lemma 4.\\


\textbf{Lemma 5:} Let $f(s)$ be a non-zero entire function, and let $s_0$ be a zero of $f(s)$. Then the multiplicity of $s_0$ is unique.

\textbf{Proof:} Let $f(s)\not\equiv 0, s\in \mathbb{C}$, be an entire function, which has a multiple zero at $s_0\in \mathbb{C}$ of multiplicity $m$.
We can write: $ f(s) = (s - s_0)^m g(s) $, where $ g(s) $ is also an entire function and $ g(s_0) \neq 0 $.

Assume for contradiction that there exists another integer $ n \neq m $ such that $ n $ is also a multiplicity of the zero $s_0$. This means we can also write: $ f(s) = (s - s_0)^n h(s) $, where $ h(s) $ is an entire function and $ h(s_0) \neq 0 $.

Since both expressions for $f(s)$ must be equal, we then obtain $ (s - s_0)^m g(s) = (s - s_0)^n h(s) $.
Without loss of generality, consider $ m > n $, then we have:
$ (s - s_0)^{m-n} g(s) = h(s) \Rightarrow h(s_0) = 0$, which is a contradiction to $h(s_0) \neq 0$.
Thus, the assumption is false, i.e., the multiplicity of a zero of any non-zero entire function is unique.

That completes the proof of Lemma 5.


\section{Foundational Theorems}
As pointed out in Ref.[2] (p.57), the non-trivial zeros of the zeta function can be enumerated in order of increasing absolute value of their imaginary parts; where zeros
whose imaginary parts have the same absolute value are arranged arbitrarily. Consequently, the same enumeration rule applies to the zeros of entire functions in the following contents.  We therefore drop the default assumption $|\beta_1|\leq|\beta_2|\leq \cdots$ as a condition hereafter for simplicity.

\textbf{Theorem 1:}
Let $F(s)$ be an entire function satisfying $F(s)=F(k-s)$, where $k >0 $ is a real constant, and admitting the following representation, which converges absolutely and uniformly on compact subsets of $\mathbb C$:
\begin{equation}
F(s)=\prod_{i=1}^{\infty}\bigl(1+\frac{(s-\alpha_i)^2}{\beta_i^2}\bigr)^{m_i},
\quad s\in\mathbb{C}
\end{equation}
Here
$m_i \geq 1$ is the multiplicity of the $i$-th quadruple of zeros $(\rho_i=\alpha_i + j\beta_i, \bar\rho_i, k-\rho_i, k-\bar\rho_i)$ of $F(s)$, $0\leq \alpha_i \leq k$, $\beta_i \neq 0$,
 $\displaystyle\sum_{i=1}^{\infty}\frac{m_i}{\beta_i^2}<\infty$.

Then
\begin{equation}
\alpha_i=\frac{k}{2}, i\in\mathbb{N}, 0<|\beta_1|<|\beta_2|<|\beta_3|<\cdots
\end{equation}

\begin{proof}
Considering \begin{equation}\nonumber\Big{(}1+\frac{(s-\alpha_i)^2}{\beta_i^2}\Big{)}^{m_{i}}=\underbrace{\Big{(}1+\frac{(s-\alpha_i)^2}{\beta_i^2}\Big{)}\cdots }_{m_i \text{ times}}\end{equation}
we have from $F(s)=F(k-s)$ and Eq.(1)
\begin{equation}
\prod_{i=1}^{\infty}\underbrace{\Big{(}1+\frac{(s-\alpha_i)^2}{\beta_i^2}\Big{)}\cdots }_{m_i \text{ times}}=\prod_{i=1}^{\infty}\underbrace{\Big{(}1+\frac{(k-s-\alpha_i)^2}{\beta_i^2}\Big{)}\cdots }_{m_i \text{ times}}
\end{equation}
where the $i^{th}$ factor appears $m_i$ times in any order.

Based on the conditions $0\leq\alpha_i\leq k$, $\sum_{i=1}^{\infty}\frac{m_i}{\beta_i^2}<\infty$ and the Weierstrass M-Test, both sides of Eq.(3) are infinite products absolutely and uniformly convergent on compact subsets of $\mathbb{C}$. Hence each side can be rearranged into any single quadratic polynomial factor times the remaining product, which stays absolutely and uniformly convergent on compact subsets of $\mathbb{C}$ and hence defines an entire function. Then by the definition of divisibility of entire functions (see Ref.[3, p.174] or Ref.[4]), Eq.(3) implies that each quadratic polynomial factor on either side divides the infinite product on the opposite side, i.e.,
 \begin{equation}
 \begin{cases}& \Big{(}1+\frac{(s-\alpha_i)^2}{\beta_i^2}\Big{)}\mid \prod_{i=1}^{\infty}\underbrace{\Big{(}1+\frac{(k-s-\alpha_i)^2}{\beta_i^2}\Big{)}\cdots }_{m_i \text{ times}}\\
 &\Big{(}1+\frac{(k-s-\alpha_i)^2}{\beta_i^2}\Big{)}\mid \prod_{i=1}^{\infty}\underbrace{\Big{(}1+\frac{(s-\alpha_i)^2}{\beta_i^2}\Big{)}\cdots }_{m_i \text{ times}}\\
 \end{cases}
 \end{equation}
 where "$\mid$" is the divisible sign, $i\in\mathbb{N}=\{1, 2, 3, \cdots\}$.

 Both polynomials $1+\frac{(s-\alpha_i)^2}{\beta_i^2}$ and $1+\frac{(k-s-\alpha_i)^2}{\beta_i^2}$ have discriminant $\Delta=-4\cdot\frac{1}{\beta_i^2}< 0$. Hence they are irreducible in $\mathbb{R}[x]$. By Lemma 2, Eq.(4) yields:
\begin{equation}\begin{cases}&\Big{(}1+\frac{(s-\alpha_i)^2}{\beta_i^2}\Big{)}\mid \Big{(}1+\frac{(k-s-\alpha_l)^2}{\beta_l^2}\Big{)}\\
 &\Big{(}1+\frac{(k-s-\alpha_i)^2}{\beta_i^2}\Big{)}\mid \Big{(}1+\frac{(s-\alpha_l)^2}{\beta_l^2}\Big{)}\\
 & i, l\in\mathbb{N}
 \end{cases}\end{equation}

For the special kind of polynomials in Eq.(5), "divisible"  means "equal", which can be verified by comparing the like terms in the following equation to get $k'=1$.
\begin{equation}\Big{(}1+\frac{(s-\alpha_i)^2}{\beta_i^2}\Big{)}=k' \Big{(}1+\frac{(k-s-\alpha_l)^2}{\beta_l^2}\Big{)}, k'\neq 0 \in \mathbb{R}
\end{equation}

Further, due to the uniqueness of the multiplicity $m_i$, the only solution to Eq.(5) is
\begin{equation}\begin{cases}&\Big{(}1+\frac{(s-\alpha_i)^2}{\beta_i^2}\Big{)}\mid \Big{(}1+\frac{(k-s-\alpha_i)^2}{\beta_i^2}\Big{)}\\
 &\Big{(}1+\frac{(k-s-\alpha_i)^2}{\beta_i^2}\Big{)}\mid \Big{(}1+\frac{(s-\alpha_i)^2}{\beta_i^2}\Big{)}
 \end{cases}, i\in\mathbb{N}\end{equation}
i.e.,
\begin{equation}1+\frac{(s-\alpha_i)^2}{\beta_i^2}=1+\frac{(k-s-\alpha_i)^2}{\beta_i^2}, i\in\mathbb{N}\tag{7a}\end{equation}
otherwise, duplicated zeros (in quadruple) with $\alpha_i + \alpha_l = k, \beta_i^2 = \beta_l^2, l \neq i$ would be generated to change $m_i$, as a result of
\begin{equation}\begin{cases}&\Big{(}1+\frac{(s-\alpha_i)^2}{\beta_i^2}\Big{)}\mid \Big{(}1+\frac{(k-s-\alpha_l)^2}{\beta_l^2}\Big{)}\\
 &\Big{(}1+\frac{(k-s-\alpha_i)^2}{\beta_i^2}\Big{)}\mid \Big{(}1+\frac{(s-\alpha_l)^2}{\beta_l^2}\Big{)}
 \end{cases}, i, l\in\mathbb{N}, i\neq l\tag{7b}\end{equation}
 i.e.,
\begin{equation}\begin{cases}&1+\frac{(s-\alpha_i)^2}{\beta_i^2}= 1+\frac{(k-s-\alpha_l)^2}{\beta_l^2}\\
 &1+\frac{(k-s-\alpha_i)^2}{\beta_i^2}= 1+\frac{(s-\alpha_l)^2}{\beta_l^2}
 \end{cases}, i, l\in\mathbb{N}, i\neq l\tag{7c}\end{equation}
By comparing the like terms in Eq.(7a), we obtain $\alpha_i = \frac{k}{2}, i\in \mathbb{N}$.
Further, to ensure the uniqueness of $m_i$ while $\alpha_i = \frac{k}{2}$, we need to restrict the $\beta_i$ values to be distinct, i.e., $0 < |\beta_1| < |\beta_2| < |\beta_3| < \cdots$.\\
That completes the proof of Theorem 1.
\end{proof}

\textbf{Remark:} Theorem 1 (with $k=1$) applies to the study of the zero distribution of the Riemann \(\xi\)-function. Let \(\rho_i=\alpha_i+j\beta_i\), with \(\beta_i\neq 0, 0<\alpha_i<1\), and let \(m_i\) be the multiplicity of the $i$-th quadruple of zeros $(\rho_i, \bar\rho_i, 1-\rho_i, 1-\bar\rho_i)$. The Hadamard product of \(\xi(s)\), grouped into conjugate pairs, can be written as
\[
\begin{aligned}
\xi(s)
&=\xi(0)\prod_{i=1}^{\infty}
\left[\left(1-\frac{s}{\rho_i}\right)
\left(1-\frac{s}{\bar\rho_i}\right)\right]^{m_i}\\
&=\xi(0)\prod_{i=1}^{\infty}
\left[
\frac{\beta_i^2}{\alpha_i^2+\beta_i^2}
\left(1+\frac{(s-\alpha_i)^2}{\beta_i^2}\right)
\right]^{m_i}
\end{aligned}
\]
Since \(\xi(s)=\xi(1-s)\), we obtain
\[
\prod_{i=1}^{\infty}
\left(1+\frac{(s-\alpha_i)^2}{\beta_i^2}\right)^{m_i}
=
\prod_{i=1}^{\infty}
\left(1+\frac{(1-s-\alpha_i)^2}{\beta_i^2}\right)^{m_i}
\]
where \(\sum_{i=1}^{\infty}\frac{m_i}{\beta_i^2}<\infty\).\\

\textbf{Theorem 2:} Let $G(s)$ be an entire function satisfying $G(s)=G(k-s)$, where $k >0 $ is a real constant, and admitting the following representation, which converges absolutely and uniformly on compact subsets of $\mathbb C$:
\begin{equation}\begin{split}
G(s)
&=\prod_{i=1}^{\infty}\Big{(}1-\frac{s}{\rho_i}\Big{)}^{m_{i}}\Big{(}1-\frac{s}{\bar{\rho}_i}\Big{)}^{m_{i}}\\
&=\prod_{i=1}^{\infty}\Big{(}\frac{\beta_i^2}{\alpha_i^2 + \beta_i^2}+\frac{(s-\alpha_i)^2}{\alpha_i^2+\beta_i^2}\Big{)}^{m_{i}}\\
&=\prod_{i=1}^{\infty} \Big{(}\frac{\beta_i^2}{\alpha_i^2 + \beta_i^2}\Big{)}^{m_i}\prod_{i=1}^{\infty} \Big{(}1+\frac{(s-\alpha_i)^2}{\beta_i^2}\Big{)}^{m_i}, s \in \mathbb{C}
\end{split}\end{equation}
where $\rho_i=\alpha_i+j\beta_i$ and $\bar\rho_i=\alpha_i-j\beta_i$ ($\beta_i\neq 0$)
      are the complex-conjugate zeros of $G(s)$; $m_i \geq 1$ is the multiplicity of the $i$-th quadruple of zeros $(\rho_i, \bar\rho_i, k-\rho_i, k-\bar\rho_i)$ of $G(s)$ with $0 \leq \alpha_i \leq k$, $\sum_{i=1}^{\infty}\frac{m_i}{|\rho_i|^2}<\infty$.

Then
\begin{equation}
 \alpha_i=\frac{k}{2}, i\in\mathbb{N}, 0<|\beta_1|<|\beta_2|<|\beta_3|<\cdots
\end{equation}

\begin{proof}
According to Lemma 3, we have
\begin{equation}\begin{split}
&G(s)=G(k-s)\\
&\Leftrightarrow\\
&\prod_{i=1}^{\infty} \Big{(}\frac{\beta_i^2}{\alpha_i^2 + \beta_i^2}\Big{)}^{m_i}\prod_{i=1}^{\infty} \Big{(}1+\frac{(s-\alpha_i)^2}{\beta_i^2}\Big{)}^{m_i}=\prod_{i=1}^{\infty} \Big{(}\frac{\beta_i^2}{\alpha_i^2 + \beta_i^2}\Big{)}^{m_i}\prod_{i=1}^{\infty} \Big{(}1+\frac{(k-s-\alpha_i)^2}{\beta_i^2}\Big{)}^{m_i}\\
&\Leftrightarrow\\
&\prod_{i=1}^{\infty} \Big{(}1+\frac{(s-\alpha_i)^2}{\beta_i^2}\Big{)}^{m_i}=\prod_{i=1}^{\infty} \Big{(}1+\frac{(k-s-\alpha_i)^2}{\beta_i^2}\Big{)}^{m_i}\\
&\Leftrightarrow\\
&F(s)=F(k-s)\\
\end{split}\end{equation}

Next, the conditions of Theorem 2 match those of Theorem 1, considering
$$0 \leq \alpha_i \leq k, \displaystyle\sum_{i=1}^{\infty}\frac{m_i}{|\rho_i|^2}<\infty\Rightarrow 0\leq \alpha_i \leq k, \displaystyle\sum_{i=1}^{\infty}\frac{m_i}{\beta_i^2}<\infty$$
Then by Theorem 1, we have
\begin{equation}
 \alpha_i=\frac{k}{2}, i\in\mathbb{N}, 0<|\beta_1|<|\beta_2|<|\beta_3|<\cdots
\end{equation}
That completes the proof of Theorem 2.
\end{proof}

 In what follows, Theorems 3 and 4 are concerned with the completed $L$- functions with a self-dual mathematical object $\lambda$. Moreover, these two theorems provide a foundation for studying the distribution of zeros of the completed $L$- functions with a non-self-dual mathematical object $\lambda$.\\

\textbf{Theorem 3:} Let $\Lambda(\lambda, s)$ be an entire function satisfying $\Lambda(\lambda, s)=\varepsilon(\lambda)\Lambda(\lambda, k-s)$, where $k >0 $ is a real constant, and admitting the Hadamard product representation:
\begin{equation}
\Lambda(\lambda, s)=e^{A(\lambda)+B(\lambda)s}\prod_{i=1}^{\infty}\Big{(}1-\frac{s}{\rho_i}\Big{)}e^{\frac{s}{\rho_i}}, s \in \mathbb{C}
\end{equation}
 Here $\lambda$ denotes a mathematical object, $\varepsilon(\lambda)$ is a complex number of absolute value 1, $\rho_i=\alpha_i+j\beta_i$ and $\bar\rho_i=\alpha_i-j\beta_i$ ($\beta_i\neq 0$)
      are the complex-conjugate zeros of $\Lambda(\lambda, s)$; $m_i \geq 1$ is the multiplicity of the $i$-th quadruple of zeros $(\rho_i, \bar\rho_i, k-\rho_i, k-\bar\rho_i)$ of $\Lambda(\lambda, s)$ with $0 \leq \alpha_i \leq k$, $\sum_{i=1}^{\infty}\frac{m_i}{|\rho_i|^2}<\infty$.

Then
\begin{equation}
 \alpha_i=\frac{k}{2}, i\in\mathbb{N}, 0<|\beta_1|<|\beta_2|<|\beta_3|<\cdots
\end{equation}

\textbf{Remark:} For more details about $\varepsilon(\lambda)$, see Ref.[5] (p.94).

\begin{proof}
First, we have
\begin{equation}\begin{split}
\Lambda(\lambda, s)&=e^{A(\lambda)+B(\lambda)s}\prod_{i=1}^{\infty}\Big{(}1-\frac{s}{\rho_i}\Big{)}e^{\frac{s}{\rho_i}}\\
&=e^{A(\lambda)+B(\lambda)s}\prod_{i=1}^{\infty}\Big{(}1-\frac{s}{\rho_i}\Big{)}\Big{(}1-\frac{s}{\bar{\rho}_i}\Big{)}e^{\frac{2\alpha_i s}{\alpha_i^2+\beta_i^2}}\\
&=e^{A(\lambda)+B(\lambda)s}e^{s\cdot\sum_{i=1}^{\infty}\frac{2\alpha_i}{\alpha_i^2+\beta_i^2}}\prod_{i=1}^{\infty}\Big{(}1-\frac{s}{\rho_i}\Big{)}\Big{(}1-\frac{s}{\bar{\rho}_i}\Big{)}\\
\end{split}\end{equation}
Noting that $\sum_{i=1}^{\infty}\frac{2\alpha_i}{\alpha_i^2+\beta_i^2}\leq 2k\cdot\sum_{i=1}^{\infty}\frac{1}{|\rho_i|^2}<\infty$, then we have $\sum_{i=1}^{\infty}\frac{2\alpha_i}{\alpha_i^2+\beta_i^2}=c, c\in \mathbb{R}$.
Further, taking into account the multiplicity $m_i\geq 1$ of each zero, we rewrite the product as
\begin{equation}\begin{split}
\Lambda(\lambda, s) &=e^{A(\lambda)+[B(\lambda)+c]s}\prod_{i=1}^{\infty}\Big{(}1-\frac{s}{\rho_i}\Big{)}^{m_i}\Big{(}1-\frac{s}{\bar{\rho}_i}\Big{)}^{m_i}\\
&=e^{A(\lambda)+[B(\lambda)+c]s} G(s)
\end{split}\end{equation}
Accordingly
\begin{equation}\begin{split}
\Lambda(\lambda,k-s) &=e^{A(\lambda)+[B(\lambda)+c](k-s)}\prod_{i=1}^{\infty}\Big{(}1-\frac{k-s}{\rho_i}\Big{)}^{m_i}\Big{(}1-\frac{k-s}{\bar{\rho}_i}\Big{)}^{m_i}\\
&=e^{A(\lambda)+[B(\lambda)+c](k-s)} G(k-s)
\end{split}\end{equation}
 Because $e^{A(\lambda)+[B(\lambda)+c]s}$ and $\varepsilon(\lambda)e^{A(\lambda)+[B(\lambda)+c](k-s)}$ are units in the ring of entire functions, they have no zeros and therefore do not affect the complex zeros related divisibility in the functional equation $\Lambda(\lambda, s)=\varepsilon(\lambda)\Lambda(\lambda, k-s)$. We thus conclude that Theorem 3 holds by Theorem 2, whose detailed conditions match those of Theorem 3.

In addition to Eq.(13), we still have (by canceling the complex zeros related polynomial factors on both sides of functional equation: $\Lambda(\lambda, s)=\varepsilon(\lambda)\Lambda(\lambda, k-s)$, since we already have $\alpha_i=\frac{k}{2}$)
\begin{equation}\Lambda(\lambda, s)=\varepsilon(\lambda)\Lambda(\lambda, k-s)\Longrightarrow e^{A(\lambda)+[B(\lambda)+c]s} =\varepsilon(\lambda)e^{A(\lambda)+[B(\lambda)+c](k-s)}\end{equation}

That completes the proof of Theorem 3.

\end{proof}

In the following Theorem 4, we suppose that the completed $L$-functions possess  both real and complex zeros, denoted by $\rho \in \mathcal{Z}_{r}$ ($\Im(\rho)=0$) and $\rho \in \mathcal{Z}_{c}$ ($\Im(\rho)\neq 0$), respectively. When these two zero sets have no common elements, we express their disjointness by: $\mathcal{Z}_{r} \cap \mathcal{Z}_{c} = \varnothing$, which is automatically satisfied by our definitions of $\mathcal{Z}_{r}$  and $\mathcal{Z}_{c}$ because $\mathcal{Z}_{r}$ and $\mathcal{Z}_{c}$ are mutually exclusive sets, i.e., if $\rho \in \mathcal{Z}_{r}$, then $\Im(\rho)=0$; if $\rho \in \mathcal{Z}_{c}$, then $\Im(\rho)\neq 0$.

The reason we need to consider this case is that, so far, we cannot rule out the existence of exceptional zeros (or Landau-Siegel zeros), although their numbers are very limited even if they do exist.

Denote the set of real zeros as
\begin{equation}
\mathcal{Z}_{r} = \left\{ a_n \in \mathbb{R} \mid 0 < a_n < k, \, n=1,2,\ldots,N \right\}
\end{equation}
where $ N $ is a finite natural number. This finiteness follows from the Identity Theorem, which implies that any non-zero entire function cannot have infinitely many zeros in a bounded region.

\textbf{Theorem 4:} Let $\Lambda(\lambda, s)$ be an entire function satisfying $\Lambda(\lambda, s)=\varepsilon(\lambda)\Lambda(\lambda, k-s)$,
where $k >0 $ is a real constant, and admitting the
Hadamard product representation:
\begin{equation}\begin{split}
\Lambda(\lambda, s)
&=e^{A(\lambda)+B(\lambda)s}\prod_{\rho}\Big{(}1-\frac{s}{\rho}\Big{)}e^{\frac{s}{\rho}}, s \in\mathbb{C}\\
&=e^{A(\lambda)+B(\lambda)s}\prod_{\rho \in \mathcal{Z}_{r}}\Big{(}1-\frac{s}{\rho}\Big{)}e^{\frac{s}{\rho}}\prod_{\rho \in \mathcal{Z}_{c}}\Big{(}1-\frac{s}{\rho}\Big{)}e^{\frac{s}{\rho}}\\
&=e^{A(\lambda)+[B(\lambda)+c]s}\prod_{\rho \in \mathcal{Z}_r}\left(1-\frac{s}{\rho}\right)e^{s/\rho} \prod_{i=1}^{\infty}\left(1-\frac{s}{\rho_i}\right)\left(1-\frac{s}{\bar{\rho}_i}\right)\\
&=e^{A(\lambda)+[B(\lambda)+c]s}\prod_{\rho \in \mathcal{Z}_r}\left(1-\frac{s}{\rho}\right)e^{s/\rho} \prod_{i=1}^{\infty}\left(1-\frac{s}{\rho_i}\right)^{m_i}\left(1-\frac{s}{\bar{\rho}_i}\right)^{m_i}
\end{split}\end{equation}

 Here $\lambda$ denotes a mathematical object, $\varepsilon(\lambda)$ is a complex number of absolute value 1, $c=\sum_{i=1}^{\infty}\frac{2\alpha_i}{\alpha_i^2+\beta_i^2}<\infty$.
\begin{itemize}
\item $\rho_i=\alpha_i+j\beta_i$ and $\bar\rho_i=\alpha_i-j\beta_i$ ($\beta_i \neq 0$)
      are the complex-conjugate zeros of $\Lambda(\lambda, s)$, $a_n, n=1,2,\ldots,N$ are the real zeros of $\Lambda(\lambda, s)$;
\item $m_i \geq 1$ is the multiplicity of the $i$-th quadruple of zeros $(\rho_i, \bar\rho_i, k-\rho_i, k-\bar\rho_i)$ of $\Lambda(\lambda, s)$;
\item $0 \leq \alpha_i \leq k$, $0 < a_n < k$, $\sum_{i=1}^{\infty}\frac{m_i}{|\rho_i|^2}<\infty$;
\item $\mathcal{Z}_{r} \cap \mathcal{Z}_{c} = \varnothing.$
\end{itemize}

Then
\begin{equation}
\begin{cases}\alpha_i=\frac{k}{2}, i\in\mathbb{N}, 0<|\beta_1|<|\beta_2|<|\beta_3|<\cdots\\ a_n=\frac{k}{2}, n=1\\
\end{cases}\end{equation}

i.e., all the zeros (both real and  complex) of $\Lambda(\lambda, s)$ in the critical strip $0 \leq \Re(s) \leq k$ lie on the critical line $\Re(s)=\frac{k}{2}$.\\

\textbf{Remark:} Since the number of real zeros is finite and the complex zeros occur in conjugate pairs and satisfy $\sum_{i=1}^{\infty} \frac{m_i}{|\rho_i|^2}<\infty, 0 \leq \alpha_i \leq k$, the decomposition
$\prod_{\rho}=\bigl(\prod_{\rho\in\mathcal{Z}_r}\bigr)\bigl(\prod_{\rho\in\mathcal{Z}_c}\bigr)$ is valid; both factors converge (the first as a finite product, the second as an absolutely and uniformly convergent infinite product).\\



\begin{proof}

By Theorem 3, to determine the distribution of the complex zeros of $\Lambda(\lambda, s)$, it suffices to show that the newly introduced factors $\prod_{\rho \in \mathcal{Z}_{r}}\Big{(}1-\frac{s}{\rho}\Big{)}e^{\frac{s}{\rho}}$  and $\prod_{\rho \in \mathcal{Z}_{r}}\Big{(}1-\frac{k-s}{\rho}\Big{)}e^{\frac{k-s}{\rho}}$ do not affect the complex zeros related divisibility in the functional equation $\Lambda(\lambda, s)=\varepsilon(\lambda)\Lambda(\lambda, k-s)$. This follows from the given condition $\mathcal{Z}_{r} \cap \mathcal{Z}_{c} = \varnothing$, which implies that $\prod_{\rho \in \mathcal{Z}_{r}}\Big{(}1-\frac{s}{\rho}\Big{)}e^{\frac{s}{\rho}}$ and $\prod_{\rho \in \mathcal{Z}_{c}}\Big{(}1-\frac{k-s}{\rho}\Big{)}e^{\frac{k-s}{\rho}}$ are co-prime, $\prod_{\rho \in \mathcal{Z}_{r}}\Big{(}1-\frac{k-s}{\rho}\Big{)}e^{\frac{k-s}{\rho}}$ and $\prod_{\rho \in \mathcal{Z}_{c}}\Big{(}1-\frac{s}{\rho}\Big{)}e^{\frac{s}{\rho}}$ are co-prime, according to Ref.[3] (p.174, p.208) or Ref.[4] (see its THEOREM 4).

Thus, considering that the conditions in Theorem 4 cover those in Theorem 3, we conclude by Theorem 3 that
\begin{equation}
 \Lambda(\lambda, s)=\varepsilon(\lambda)\Lambda(\lambda, k-s)\Rightarrow\begin{cases}&\alpha_i=\frac{k}{2}, i\in\mathbb{N}\\ & 0<|\beta_1|<|\beta_2|<|\beta_3|< \cdots\\
  \end{cases}
\end{equation}

Next, we consider the real zeros of $\Lambda(\lambda, s)$.

By canceling the complex zeros related polynomial factors on both sides of $\Lambda(\lambda, s)=\varepsilon(\lambda)\Lambda(\lambda, k-s)$, we have
\begin{equation}e^{A(\lambda)+[B(\lambda)+c]s}\prod_{n=1}^{N}\Big{(}1-\frac{s}{a_n}\Big{)}e^{\frac{s}{a_n}}
=\varepsilon(\lambda)e^{A(\lambda)+[B(\lambda)+c](k-s)}\prod_{n=1}^{N}\Big{(}1-\frac{k-s}{a_n}\Big{)}e^{\frac{k-s}{a_n}}
\end{equation}
where constant $c$ is the same as in the proof of Theorem 3.

Further Eq.(22) is equivalent to
\begin{equation}(-1)^Ne^{A(\lambda)+[B(\lambda)+c']s}\prod_{n=1}^{N}\Big{(}s-a_n\Big{)}
=\varepsilon(\lambda)e^{A(\lambda)+[B(\lambda)+c'](k-s)}\prod_{n=1}^{N}\Big{(}s-(k-a_n)\Big{)}
\end{equation}
where $c'=c+\sum_{n=1}^{N} \frac{1}{a_n}$.

Suppose the multiplicity of zero $s=a_n$ is $m_n$ ($m_n\geq 1$) that is finite and uniquely determined although unknown. Then Eq.(23) becomes
 \begin{equation}(-1)^Ne^{A(\lambda)+[B(\lambda)+c']s}\prod_{t=1}^{T}\Big{(}s-a_t\Big{)}^{m_t}
=\varepsilon(\lambda)e^{A(\lambda)+[B(\lambda)+c'](k-s)}\prod_{t=1}^{T}\Big{(}s-(k-a_t)\Big{)}^{m_t}
\end{equation}
where $\sum_{t=1}^{T} m_t = N$.

Considering $(s-a_t)$ and $(s-(k-a_t))$ are irreducible in $\mathbb{R}[x]$, then by Lemma 1 or Lemma 2, Eq.(24) means
\begin{equation}\begin{cases}&(s-a_t)\mid (s-(k-a_l)\\
 &(s-(k-a_t))\mid (s-a_l)\\
 & t, l=1, 2, 3, \cdots, T
 \end{cases}\end{equation}

The only solution to Eq.(25) is $t=l, a_t=\frac{k}{2}, t=1, \cdots, T$, otherwise the uniqueness of $m_t$ would be violated with $a_t+a_l=k, t\neq l$.
 To avoid changing the multiplicity of $m_t$ while $a_t=\frac{k}{2}$, we need to limit $T=1$. Thus we get
\begin{equation} \Lambda(\lambda, s)=\varepsilon(\lambda)\Lambda(\lambda, k-s)\Rightarrow a_t=\frac{k}{2}, t=1, m_1=N
\end{equation}

Putting Eq.(21) and Eq.(26) together, we proved Eq.(20).

In addition to Eq.(20), we still have (by canceling the real zeros related polynomial factors on both sides of Eq.(24))
\begin{equation}\Lambda(\lambda, s)=\varepsilon(\lambda)\Lambda(\lambda, k-s)\Longrightarrow (-1)^Ne^{A(\lambda)+[B(\lambda)+c']s}=\varepsilon(\lambda)e^{A(\lambda)+[B(\lambda)+c'](k-s)}
\end{equation}
That completes the proof of Theorem 4.
\end{proof}

\noindent We have the following supplementary remarks regarding Theorem 4.

\textbf{1):} If $\Lambda(\lambda, s)$ has poles/zeros at $s=0, k$ with multiplicity (order) $r\geq 1$, then $s^r(k-s)^r\Lambda(\lambda, s)$ is an entire function. This adjustment does not affect the results of Theorem 4, since we have $s^r(k-s)^r\Lambda(\lambda, s)=\varepsilon(\lambda)(k-s)^r s^r\Lambda(\lambda, k-s)\Leftrightarrow \Lambda(\lambda, s)=\varepsilon(\lambda)\Lambda(\lambda, k-s)$.

\textbf{2):} As pointed out in Ref.[5] (p.102), if $\rho_i= \alpha_i +j\beta_i$ is a zero of $\Lambda(\lambda, s)$, then $\bar{\rho}_i= \alpha_i -j\beta_i$ is a zero of $\Lambda(\overline{\lambda}, s)$, $\overline{\lambda}$ is the dual of $\lambda$. Therefore, to use Theorem 4 while $\lambda\neq\bar{\lambda}$ with symmetric functional equation $\Lambda(\lambda, s)=\varepsilon(\lambda)\Lambda(\overline\lambda, k-s)$, we need to construct a new symmetric functional equation $\Lambda(\overline{\lambda}, s) \Lambda(\lambda, s)= \varepsilon(\lambda)\varepsilon(\overline\lambda) \Lambda(\lambda, k-s)\Lambda(\overline{\lambda}, k-s)$ to ensure that the conjugate zeros appear together in the related Hadamard products.

\textbf{3):} Theorems 1-4, can not be applied to the Davenport-Heilbronn counterexample. That is because, in the proofs of Theorems 1-4, $\sum_{i=1}^{\infty}\frac{m_i}{|\beta_i|^2}<\infty, 0\leq \alpha_i \leq k$ is a necessary condition.
However, the Davenport-Heilbronn counterexample does not guarantee that all its zeros lie inside the critical strip since it possesses no Euler product (see Ref.[5] on page 102 for more details); therefore, one cannot deduce that $\sum_{i=1}^{\infty}\frac{m_i}{|\beta_i|^2}<\infty$ from $\sum_{i=1}^{\infty}\frac{m_i}{|\rho_i|^2}<\infty$ without $0\leq \alpha_i \leq k$.

%\textbf{Actually, Theorem 4 provides a unified proof framework for the Extended Riemann Hypothesis, the Generalized %Riemann Hypothesis, and the Grand Riemann Hypothesis.}

\section{Main Results}

We will make use of Theorem 4 to prove the extended Riemann Hypothesis (for Dedekind zeta functions), the generalized Riemann Hypothesis (for Dirichlet \( L \)-functions), and the grand Riemann Hypothesis (for cuspidal automorphic \( L \)-functions, and for any kind of \( L \)-functions, repectively).

In the following contents, the critical line means: $\mathfrak{R}(s)=\frac{1}{2}$, or more generally, $\mathfrak{R}(s)=\frac{k}{2}$, $k > 0$ is a real constant.

Ref.[5] (mainly in Section 5.1) provides a detailed account of the general properties of L-functions, including the unified definition (via Dirichlet series with Euler product), the gamma factor, the conductor, the root number, the completed \( L \)-function, the symmetric functional equation, the Hadamard product (as an entire function of order 1), among others.

In addition, we also have the following general property of \( L \)-functions, which was labeled Lemma 5.5 in Ref.[5] (p.101).

\textbf{Lemma 5.5 }$^{[5]}$:
Let \( L(f,s) \) be an \( L \)-function. All zeros \( \rho \) of \( \Lambda(f,s) \) are in the critical strip \( 0 \leq \sigma \leq 1 \). For any \( \epsilon > 0 \), we have
\[
\sum_{\rho \neq 0,1} |\rho|^{-1-\epsilon} < +\infty.
\]
In Lemma 5.5, \( \Lambda(f,s) \) is the completed \( L \)-function corresponding to \( L(f,s) \), $\sigma$ is the real part of \( \rho \), and $f$ is identical to $\lambda$ in this paper as a symbol representing a mathematical object (e.g., Dirichlet character, modular form, automorphic representation).

Actually, in Lemma 5.5, the critical strip can be modified from the closed interval \( 0 \leq \sigma \leq 1 \) to the open interval \( 0 < \sigma < 1 \), according to Refs.[10-12].

To adapt to the wider situation of critical strip, i.e., $0 \leq \Re(s) \leq k$, Lemma 5.5 has been extended to Theorem 9 (positioned after Theorem 8).

Another general property of \( L \)-functions is as follows.

The zeros of the completed $L$-function $\Lambda(\lambda, s)$ are precisely the non-trivial zeros of $L(\lambda, s)$. See Ref.[5] (p.96) for more details (with $f$ therein replaced by $\lambda$).

\textbf{Thus, we can discuss the non-trivial zeros of $L$-functions based on the zeros of the corresponding completed $L$-functions.}

\subsection{Dirichlet $L$-Function}
We briefly list the contents of Dirichlet $L$-functions relevant to this paper, based mainly on Ref.[5].

\textbf{Definition:}
Let $\chi$ be a Dirichlet character modulo $q$. The Dirichlet $L$-function associated with $\chi$ is defined by
\begin{equation}
L(\chi, s)
=
\sum_{n=1}^{\infty}\frac{\chi(n)}{n^s}
=
\prod_{p}
\left(1-\frac{\chi(p)}{p^s}\right)^{-1},
\qquad \Re(s)>1.
\end{equation}
where the product is taken over all primes $p$.

For the principal (trivial) character $\chi_0$ (where $\chi_0(n) = 1$ if $\gcd(n,q) = 1$ and $\chi_0(n) = 0$ otherwise), the $L$-function is related to the Riemann zeta function by:

\begin{equation}L(\chi_0, s) = \zeta(s) \prod_{p|q} \left(1 - \frac{1}{p^s}\right)\end{equation}

\textbf{Remark:} The Riemann zeta-function $\zeta(s)$ is a special case of $L(\chi, s)$ with $\chi(n) \equiv 1$ ($q = 1$).\\

\textbf{Completed L-function:}
Let $\chi$ be a primitive Dirichlet character modulo $q > 1$. The completed Dirichlet $L$-function is defined as:

\begin{equation} \Lambda(\chi, s) = q^{\frac{s}{2}}\pi^{-\frac{s}{2}} \Gamma\left(\frac{s+a}{2}\right) L( \chi, s) =\left(\frac{q}{\pi}\right)^{\frac{s}{2}} \Gamma\left(\frac{s+a}{2}\right) L( \chi, s) \end{equation}

where $ a = 0 $ if $ \chi(-1) = 1 $ (even character) and $ a = 1 $ if $ \chi(-1) = -1 $ (odd character).

\textbf{Functional Equation:}
The completed Dirichlet $L$-function satisfies the functional equation:

\begin{equation} \Lambda(\chi, s) = W(\chi) \Lambda(\overline{\chi}, 1-s) \end{equation}

where $ W(\chi) $ is the root number:
\begin{equation} W(\chi) = \frac{\tau(\chi)}{j^a \sqrt{q}} \end{equation}

with $j=\sqrt{-1}, \tau(\chi) = \sum_{n=1}^{q} \chi(n) e^{2\pi j n/q} $ being the Gauss sum associated with $ \chi $.

\textbf{Hadamard Product:}
For primitive non-principal character $\chi$, the completed $L$-function $\Lambda(\chi, s)$ is an entire function of order $1$ and has the Hadamard product:

\begin{equation}\Lambda(\chi, s) = e^{A(\chi) + B(\chi)s} \prod_{\rho} \left(1 - \frac{s}{\rho}\right) e^{s/\rho} \end{equation}

where the product is over all zeros $\rho$ of $\Lambda(\chi, s)$, and $A(\chi)$ and $B(\chi)$ are constants depending on $\chi$.

For the principal character $\chi_0$ modulo $q > 1$, Eq.(29) shows that the
non-trivial zeros of $L(\chi_0, s)$ coincide with those of $\zeta(s)$, since
the finite Euler product $\prod_{p \mid q} \left(1 - p^{-s}\right)$ does not
vanish in the critical strip; this case thus reduces to the case $q = 1$.

For $q = 1$, the trivial character $\chi_0$ satisfies $\chi_0(n) \equiv 1$,
so that $L(\chi_0, s) = \zeta(s)$ and, with $a = 0$,
\begin{equation}\nonumber
\Lambda(\chi_0, s) = \pi^{-s/2} \Gamma\!\left(\frac{s}{2}\right) \zeta(s),
\end{equation}
which carries simple poles at $s = 0$ and $s = 1$. Since the factor $s(1-s)$
is invariant under $s \mapsto 1-s$, the function
\begin{equation}\nonumber
\Lambda'(\chi_0, s) = s(1-s) \Lambda(\chi_0, s)
\end{equation}
is an entire function of order $1$, and the functional equation
$\Lambda(\chi_0, s) = \Lambda(\overline{\chi}_0, 1-s)$ with
$\overline{\chi}_0 = \chi_0$ and root number $W(\chi_0) = 1$ implies the
symmetric functional equation
\begin{equation}\nonumber
\Lambda'(\chi_0, s) = \Lambda'(\chi_0, 1-s).
\end{equation}
Hence the Hadamard product applies to $\Lambda'(\chi_0, s)$. This is
exactly the situation considered in the proof of Theorem 6, and it makes
no difference to our result; thus this (trivial) case is omitted in the
proof of Theorem 5.

Next we prove the generalized Riemann Hypothesis, noting that $0 < \Re(s) < 1$ is a subset of $0 \leq \Re(s) \leq 1$.

\textbf{Theorem 5:} The non-trivial zeros of Dirichlet $L$-functions in the critical strip $0 \leq \Re(s) \leq 1$ lie on the critical line.

\textbf{Remark:} The non-trivial zero problem of $L(\chi, s)$ for arbitrary non-principal characters reduces to that for primitive non-principal characters, so we only consider primitive non-principal characters in the following proof.

\textbf{Remark:} It suffices to prove that all the zeros of $\Lambda(\chi, s)$ in the critical strip $0 \leq \Re(s) \leq 1$ have real part $\frac{1}{2}$, i.e., all the zeros of $\Lambda(\chi, s)$ lie on the critical line.

\begin{proof} We conduct the proof in two cases.

 CASE 1: $\chi=\bar{\chi}$ (self-dual)

 It suffices to verify that the properties of $\Lambda(\chi, s)$  match the following conditions of Theorem 4 with $\lambda=\chi$, $\chi=\bar{\chi}$, $\varepsilon(\lambda)=W(\chi)$, $k=1$:

 1) Symmetric functional equation;

 2) Hadamard product expression;

 3) The complex zeros of $\Lambda(\chi, s)$ appear in pairs $(\rho_i, \bar\rho_i)$ with multiplicity $m_i\ge 1$. Furthermore $m_i$ is the multiplicity of the $i$-th quadruple of zeros ($\rho_i, \bar\rho_i, 1-\rho_i, 1-\bar\rho_i$) of $\Lambda(\chi, s)$ under the symmetric functional equation;

 4) $0 \leq \alpha_i \leq 1$, $\displaystyle\sum_{i=1}^{\infty}\frac{1}{|\rho_i|^2}<\infty$, $0 < a_n < 1, n\leq N$;

 5) $\mathcal{Z}_{r} \cap \mathcal{Z}_{c} = \varnothing$.

   Eq.(31) shows the symmetric functional equation; Hadamard product Eq.(33) is equivalent to Eq.(19) in Theorem 4 by separating all zeros into two sets $\mathcal{Z}_{r}$ and  $\mathcal{Z}_{c}$;
 $\chi=\bar{\chi}$  guarantees that the complex conjugate zeros of $\Lambda(\chi, s)$ appear in pairs $(\rho_i, \bar\rho_i)$ with multiplicity $m_i\ge 1$, and then $m_i$ is the multiplicity of the $i$-th quadruple of zeros ($\rho_i, \bar\rho_i, 1-\rho_i, 1-\bar\rho_i$) by Eq.(31); The condition $0 \leq \alpha_i \leq 1$, $\sum_{i=1}^{\infty}\frac{1}{|\rho_i|^2}<\infty$, $0 < a_n < 1$ can be assured by Lemma 5.5, considering that $\sum_{i=1}^{\infty}\frac{1}{|\rho_i|^2}$  is a subseries of $\sum_{\rho\neq 0, 1}\frac{1}{|\rho|^2}$; The condition
 $\mathcal{Z}_{r} \cap \mathcal{Z}_{c} = \varnothing$ holds because $\mathcal{Z}_{r}$ and $\mathcal{Z}_{c}$ are mutually exclusive sets, i.e., if $\rho \in \mathcal{Z}_{r}$, then $\Im(\rho)=0$; if $\rho \in \mathcal{Z}_{c}$, then $\Im(\rho)\neq 0$.

Therefore, by Theorem 4 with $\lambda=\chi$, $\chi=\bar{\chi}$, $\varepsilon(\lambda)=W(\chi)$, $k=1$, we know that all zeros (real, if any, and complex) of $\Lambda(\chi, s)$ in the critical strip $0 \leq \Re(s) \leq 1$ lie on the critical line.

 CASE 2: $\chi\neq\bar{\chi}$ (non-self-dual)

 In this case, the complex conjugate zeros do not appear together in Eq.(33), because if $\rho$ is a zero of $\Lambda(\chi, s)$, then $\bar{\rho}$ is a zero of $\Lambda(\bar{\chi}, s)$. To address this issue, we need to construct a new entire function $\Lambda(\chi, s)\Lambda(\bar{\chi}, s)$ with corresponding symmetric functional equation and Hadamard product expression.

 First, we extend Eq.(31) to another form, i.e.,
 \begin{equation} \Lambda(\bar{\chi}, s) = W(\bar{\chi}) \Lambda(\chi, 1-s) \end{equation}
 Combining Eq.(34) with Eq.(31), we get a new symmetric functional equation
 \begin{equation} \Lambda(\bar{\chi}, s) \Lambda(\chi, s)= W(\bar{\chi})W(\chi) \Lambda(\chi, 1-s)\Lambda(\bar{\chi}, 1-s) \end{equation}

 The product $\Lambda(\bar{\chi}, s)\Lambda(\chi, s)$ is an entire function of finite order $\leq 1$, as this holds for the product of any two entire functions of order 1.
 Moreover, the zero set of $\Lambda(\bar{\chi}, s)\Lambda(\chi, s)$—being the union of the zeros of the individual factors—possesses the conjugate zeros paring property.

 Then all complex zeros of $\Lambda(\bar{\chi}, s)\Lambda(\chi, s)$  come in conjugate pairs $(\rho_i, \bar\rho_i)$ with multiplicity $m_i$; furthermore, by Eq.(35), $m_i$ is the multiplicity of the $i$-th quadruple of zeros ($\rho_i, \bar\rho_i, 1-\rho_i, 1-\bar\rho_i$). In addition, Lemma 5.5 applies to $\Lambda(\bar{\chi}, s)\Lambda(\chi, s)$.

 To be specific, based on Eq.(33), we have the following Hadamard product expression for $\Lambda(\chi, s)\Lambda(\bar{\chi}, s)$
\begin{equation}\begin{split}
\Lambda(\chi, s) \Lambda(\bar{\chi}, s)
&= \left[ e^{A(\chi) + B(\chi)s} \prod_{\rho \in \mathcal{Z}_r} \left(1 - \frac{s}{\rho}\right)e^{s/\rho} \prod_{\rho \in \mathcal{Z}_c} \left(1 - \frac{s}{\rho}\right)e^{s/\rho} \right] \\
&\quad \cdot \left[ e^{A(\bar\chi) + B(\bar\chi)s} \prod_{\bar\rho \in \mathcal{Z}_r} \left(1 - \frac{s}{\bar\rho}\right)e^{s/\bar\rho} \prod_{\bar\rho \in \mathcal{Z}_c} \left(1 - \frac{s}{\bar\rho}\right)e^{s/\bar\rho} \right] \\
&= e^{A(\chi) + A(\bar\chi) + [B(\chi)+B(\bar\chi)+c]s} \prod_{\rho \in \mathcal{Z}_r}\left(1-\frac{s}{\rho}\right)e^{s/\rho} \prod_{i=1}^{\infty}\left(1-\frac{s}{\rho_i}\right)\left(1-\frac{s}{\bar{\rho}_i}\right)
\end{split}\end{equation}

 where $c=\sum_{i=1}^{\infty}\frac{2\alpha_i}{\alpha_i^2+\beta_i^2}$.

 The other conditions required by Theorem 4 hold for the same reasons as in CASE 1.

Therefore, according to Theorem 4, all zeros (real, if any, and complex) of $\Lambda(\chi, s) \Lambda(\bar{\chi}, s)$, and consequently of $\Lambda(\chi, s)$,  in the critical strip $0 \leq \Re(s) \leq 1$, lie on the critical line for $\chi\neq\bar{\chi}$.

 Combining CASE 1 and CASE 2, we conclude that Theorem 5 holds as a specific case of Theorem 4 with $\lambda=\chi$,  $\varepsilon(\lambda)=W(\chi)$, $k=1$.
\end{proof}

\textbf{Remark:} According to Theorem 5: All non-trivial zeros (both real and complex) of Dirichlet $L$-functions, in the critical strip $0\leq \Re(s) \leq 1$, lie on the critical line $\Re(s)=\frac{1}{2}$.  This implies the non-existence of Landau-Siegel zeros. Hence, the Landau-Siegel zeros conjecture is justified.

\subsection{ Dedekind Zeta Function}
We briefly list the contents of Dedekind zeta functions relevant to this paper, based mainly on Ref.[5].

\textbf{Definition:}
For a number field $K/\mathbb{Q}$ of degree
$d=[K:\mathbb{Q}]$ with ring of integers
$\mathcal{O}_K$, the Dedekind zeta function is defined by
\begin{equation}
\zeta_K(s)
=
\sum_{\mathfrak{a}\neq 0}
\frac{1}{(N\mathfrak{a})^s}
=
\prod_{\mathfrak{p}}
\left(1-\frac{1}{(N\mathfrak{p})^s}\right)^{-1},
\qquad \Re(s)>1.
\end{equation}
where the sum is over all nonzero integral ideals
$\mathfrak{a}$ of $\mathcal{O}_K$, the product is over
all nonzero prime ideals $\mathfrak{p}$ of
$\mathcal{O}_K$, and
$N\mathfrak{a}=|\mathcal{O}_K/\mathfrak{a}|$
denotes the norm of the ideal
$\mathfrak{a}$.

\textbf{Remark:} The Riemann zeta-function $\zeta(s)$ is a special case of $\zeta_K(s)$ with $K=\mathbb{Q}$, where $\mathbb{Q}$ denotes the field of rational numbers.\\

\textbf{Completed Zeta Function:}
The completed Dedekind zeta function is defined as:

\begin{equation}\Lambda_K(s) = |D_K|^{s/2} \pi^{-\frac{ds}{2}}\Gamma\left(\frac{s}{2}\right)^{r_1+r_2}
\Gamma\left(\frac{s+1}{2}\right)^{r_2} \zeta_K(s)\end{equation}

where $D_K$ is the discriminant of $K$, $r_1$ is the number of real embeddings of $K$, $r_2$ is the number of pairs of complex embeddings of $K$, $d=r_1+2r_2$.


\textbf{Functional Equation:}
The completed Dedekind zeta function satisfies the functional equation

\begin{equation}\Lambda_K(s) = \Lambda_K(1-s)\end{equation}


\textbf{Hadamard Product:}
The completed Dedekind zeta function has a simple pole at $s = 1$. The function $s(1-s)\Lambda_K(s)$ is an entire function of order 1 and has the Hadamard product:

\begin{equation}s(1-s)\Lambda_K(s)= e^{A_K + B_K s} \prod_{\rho \neq 0, 1} \left(1 - \frac{s}{\rho}\right) e^{s/\rho}\end{equation}

where the product runs over all zeros $\rho$ of $\Lambda_K(s)$, and $A_K$ and $B_K$ are constants depending on $K$.\\

For more details of Dedekind zeta functions, please refer to Ref.[5] (Section 5.1 and Section 5.10) and Ref.[6] (Section 10.5.1).  \\

Next, we prove the extended Riemann Hypothesis, noting that $0 < \Re(s) < 1$ is a subset of $0 \leq \Re(s) \leq 1$.

\textbf{Theorem 6:} The non-trivial zeros of Dedekind zeta functions in the critical strip $0 \leq \Re(s) \leq 1$ lie on the critical line.

\textbf{Remark:} It suffices to prove that all the zeros of $\Lambda_K(s)$ in the critical strip $0 \leq \Re(s) \leq 1$ have real part $\frac{1}{2}$, i.e., all the zeros of $\Lambda_K(s)$ lie on the critical line.

\begin{proof}
Let $\Lambda_K'(s)=s(1-s)\Lambda_K(s)$. Then $\Lambda_K'(s)$ is an entire function of order 1. And that, the zeros of $\Lambda_K(s)$ coincide with the zeros of $\Lambda_K'(s)$. Then, it suffices to show that the properties of $\Lambda_K'(s)$ match the conditions of Theorem 4 with $\lambda = \bar\lambda = K$, $\varepsilon(\lambda)=1$, $k=1$.

 Actually, symmetric functional equation Eq.(39) guarantees that
\begin{equation}
\Lambda_K'(s)=s(1-s)\Lambda_K(s) = (1-s)s\Lambda_K(1-s)=\Lambda_K'(1-s)
\end{equation}

Recalling Theorem 4, $\Lambda_K'(s)$ here corresponds to $\Lambda(\lambda, s)$, and Eq.(40) corresponds to Eq.(19).

Furthermore, the complex conjugate zeros of $\Lambda_k'(s)$ appear in pairs $(\rho_i, \bar\rho_i)$ with multiplicity $m_i\ge 1$, since it is self-dual as $\Lambda_K(s)$ is, and then $m_i$ is the multiplicity of the $i$-th quadruple of zeros ($\rho_i, \bar\rho_i, 1-\rho_i, 1-\bar\rho_i$) by Eq.(41); The other conditions required by Theorem 4 hold for the same reasons as in CASE 1 in the proof of Theorem 5.

Therefore, by Theorem 4, we know that all zeros (real, if any, and complex) of $\Lambda_K'(s)$  (thus of $\Lambda_K(s)$) in the critical strip $0 \leq \Re(s) \leq 1$ lie on the critical line. Thus, Theorem 6 holds as a specific case of Theorem 4 with $\lambda = \bar\lambda = K$, $\varepsilon(\lambda)=1$, $k=1$.
\end{proof}



\subsection{Automorphic $L$-Function}

We briefly list the contents of automorphic $L$-functions relevant to this paper, based mainly on Ref.[5].

\textbf{Definition:} Let $\pi \simeq \otimes'_v \pi_v$ be a cuspidal automorphic representation of $\mathrm{GL}_m(\mathbb{A}_{\mathbb{Q}})$ with unitary central character, normalized such that the central character is
trivial on the diagonally embedded copy of the positive reals, where $m\ge 1$ is an integer, $\simeq$ denotes isomorphism of representations, $\otimes'_v$ denotes the restricted tensor product over all places $v$ of $\mathbb{Q}$ (including the archimedean place $v=\infty$), and $\mathbb{A}_{\mathbb{Q}}$ denotes the adele ring of $\mathbb{Q}$. In addition, $Q_\pi$ is a positive integer, the arithmetic conductor of $\pi$.

The standard $L$-function associated to $\pi$ is defined by
\begin{equation}
L(\pi, s)=\prod_{p<\infty} L_p(\pi, s)=\sum_{n=1}^{\infty} \frac{\lambda_{\pi}(n)}{n^s}, \quad\quad\Re(s)>1.
\end{equation}
where $L_p(\pi, s)$ is the local $L$-factor at the prime $p$. For each (finite) prime $p$, the
inverse of the local factor $L_p(\pi, s)$ is a polynomial in $p^{-s}$ of degree
$\leq m$, i.e.
\begin{equation}
L_p(\pi, s)=\prod_{i=1}^m \left(1-\alpha_{i,\pi}(p)p^{-s}\right)^{-1}
\end{equation}
where $\alpha_{i, \pi}(p)$ are suitable complex numbers,
$\alpha_{i, \pi}(p) \neq 0$ for all $i$ whenever $p \nmid Q_{\pi}$ ($\pi_p$ is unramified), and it might
be the case that $\alpha_{i, \pi}(p) = 0$ for some $i$ when $p\mid Q_{\pi}$ ($\pi_p$ is ramified). Please refer to Ref.[9] for more details of coefficients $\alpha_{i, \pi}(p)$ and $\lambda_{\pi}(n)$.

At the archimedean place of $\mathbb{Q}$, there are $m$ complex Langlands
parameters $\mu_{i, \pi}$, based on which the archimedean
factor $L_{\infty}(s, \pi)$ is defined by the following product of gamma factors
\begin{equation}
L_{\infty}(s, \pi) = \pi^{-\frac{ms}{2}} \prod_{i = 1}^{m}
\Gamma\!\left( \frac{s + \mu_{i, \pi}}{2} \right). \notag
\end{equation}

\textbf{Remark:} The Riemann zeta-function $\zeta(s)$ is a special case of $L(\pi, s)$ with $m=1, \pi \simeq \mathbf{1}$ (the trivial representation). \\

\textbf{Completed $L$-function:}
The completed automorphic $L$-function is defined by
\begin{equation}
\Lambda(\pi, s)=Q_\pi^{\,s/2}L_\infty(\pi, s)L(\pi, s)
\end{equation}

\textbf{Functional Equation:}
The completed automorphic $L$-function satisfies the functional equation
\begin{equation}
\Lambda(\pi, s)=\varepsilon(\pi)\,\Lambda(\tilde{\pi}, 1-s)
\end{equation}
where $\tilde{\pi}$ is the contragredient representation of $\pi$, $\varepsilon(\pi)$ is a complex number of absolute value $1$, i.e., $|\varepsilon(\pi)|=1$,  also known as the root number of the corresponding $L$-function.

\textbf{Hadamard Product:}
For cuspidal $\pi$ on $\mathrm{GL}_m$ with $m\geq 2$, $\Lambda(\pi, s)$ is an entire function of order $1$, and hence admits the Hadamard product expansion
\begin{equation}
\Lambda(\pi, s)=e^{A(\pi)+B(\pi)s}\prod_{\rho}\left(1-\frac{s}{\rho}\right)e^{s/\rho}
\end{equation}
where the product runs over all zeros $\rho$ of $\Lambda(\pi, s)$, and $A(\pi)$ and $B(\pi)$ are constants depending on $\pi$.

For cuspidal $\pi$ on $\mathrm{GL}_1$, $s^r(1-s)^r\Lambda(\pi, s)$, rather than $\Lambda(\pi, s)$ itself, is an entire function, $r\geq 1$ is the multiplicity (order) of pole of $\Lambda(\pi, s)$ at $s=1$. This situation can be processed  as shown in the proof of Theorem 6, it doesn't affect the final conclusion. Thus we ignore this situation in the proof of Theorem 7.

For more details of automorphic $L$-functions, please refer to Refs.[5][7][8].

The grand Riemann Hypothesis (GRH) has different versions of statements. To be specific, we adopt the following two versions.

GRH (version 1): For any cuspidal representation $\pi$, the non-trivial zeros of the $L$-function $L(\pi, s)$ in the critical strip $0 < \Re(s) < 1$ lie on the critical line $\Re(s)=\frac{1}{2}$.


GRH (version 2): The non-trivial zeros of any $L$-function in the critical strip $0 < \Re(s) < 1$ lie on the critical line $\Re(s)=\frac{1}{2}$.

GRH (version 1) corresponds to Theorem 7, GRH (version 2) corresponds to Theorem 8.

\textbf{Theorem 7:} The non-trivial zeros of cuspidal automorphic $L$-Functions in the critical strip $0 \leq \Re(s) \leq 1$ lie on the critical line.

\textbf{Remark:} It suffices to prove that all the zeros of $\Lambda(\pi, s)$ in the critical strip $0 \leq \Re(s) \leq 1$ have real part $\frac{1}{2}$, i.e., all the zeros of $\Lambda(\pi, s)$ lie on the critical line.

\begin{proof} We conduct the proof in two cases.

CASE 1: $\pi \simeq \tilde{\pi}$ (self-dual, abbreviated as $\pi = \tilde{\pi}$ )

It suffices to show that the properties of $\Lambda(\pi, s)$ match the conditions of Theorem 4 with $\lambda=\pi= \tilde{\pi}=\bar\lambda, k=1$. Eq.(45) shows the symmetric functional equation; Hadamard product Eq.(46) is equivalent to Eq.(19) in Theorem 4 by separating all zeros into two sets $\mathcal{Z}_{r}$ and  $\mathcal{Z}_{c}$;
  $\pi=\tilde{\pi}$ guarantees that the complex conjugate zeros of $\Lambda(\pi, s)$ appear in pairs $(\rho_i, \bar\rho_i)$ with multiplicity $m_i \ge 1$, and then $m_i$ is the multiplicity of the $i$-th quadruple of zeros ($\rho_i, \bar\rho_i, 1-\rho_i, 1-\bar\rho_i$) by Eq.(45); The condition $0 \leq \alpha_i \leq 1$, $\sum_{i=1}^{\infty}\frac{1}{|\rho_i|^2}<\infty$, $0 < a_n < 1$ can be assured by Lemma 5.5; The condition
 $\mathcal{Z}_{r} \cap \mathcal{Z}_{c} = \varnothing$ holds as explained in the proof of Theorem 5.

 Therefore, by Theorem 4, we know that all zeros (real, if any, and complex) of $\Lambda(\pi, s)$ in the critical strip $0 \leq \Re(s) \leq 1$ lie on the critical line.


CASE 2: $\pi \not\simeq \tilde{\pi}$ (non-self-dual)

In this case, the complex conjugate zeros do not appear together in Eq.(46), because if $\rho$ is a zero of $\Lambda(\pi, s)$, then $\bar{\rho}$ is a zero of $\Lambda(\tilde{\pi}, s)$. To address this issue, we need to construct a new entire function $\Lambda(\tilde{\pi}, s) \Lambda(\pi, s)$ with corresponding symmetric functional equation and Hadamard product expression.

We first extend Eq.(45) to another form, i.e.,
 \begin{equation} \Lambda(\tilde{\pi}, s) = \varepsilon(\tilde{\pi}) \Lambda(\pi, 1-s) \end{equation}
 Combining Eq.(47) with Eq.(45), we get a new symmetric functional equation
 \begin{equation} \Lambda(\tilde{\pi}, s) \Lambda(\pi, s)= \varepsilon(\pi)\varepsilon(\tilde{\pi}) \Lambda(\pi, 1-s)\Lambda(\tilde{\pi}, 1-s) \end{equation}
 Both sides of Eq.(48) are the products of entire functions of order 1, thus they are still entire functions of order $\leq 1$, as this holds for the product of any two entire functions of order 1.
 Moreover, the zero set of $\Lambda(\tilde{\pi}, s) \Lambda(\pi, s)$—being the union of the zeros of the individual factors—possesses the conjugate zeros paring property.

 Then all complex zeros of $\Lambda(\tilde{\pi}, s) \Lambda(\pi, s)$  come in conjugate pairs $(\rho_i, \bar\rho_i)$ with multiplicity $m_i$; furthermore, by Eq.(48), $m_i$ is the multiplicity of the $i$-th quadruple of zeros ($\rho_i, \bar\rho_i, 1-\rho_i, 1-\bar\rho_i$). In addition, Lemma 5.5 applies to $\Lambda(\tilde{\pi}, s) \Lambda(\pi, s)$.

 Based on Eq.(46), we have the following Hadamard product expression for $\Lambda(\tilde{\pi}, s) \Lambda(\pi, s)$
 \begin{equation}\Lambda(\tilde{\pi}, s) \Lambda(\pi, s) = e^{A(\pi) + A(\tilde{\pi})+[B(\pi)+B(\tilde{\pi})+c]s}\prod_{\rho \in \mathcal{Z}_{r}}\Big{(}1-\frac{s}{\rho}\Big{)}e^{\frac{s}{\rho}} \prod_{i=1}^{\infty} \left(1 - \frac{s}{\rho_i}\right)\left(1 - \frac{s}{\bar{\rho}_i}\right)  \end{equation}
 where $c=\sum_{i=1}^{\infty}\frac{2\alpha_i}{\alpha_i^2+\beta_i^2}$.

 The other conditions required by Theorem 4 hold for the same reasons as in CASE 1.

 Therefore, by Theorem 4, all zeros (real, if any, and complex) of $\Lambda(\tilde{\pi}, s) \Lambda(\pi, s)$, and consequently of $\Lambda(\pi, s)$, in the critical strip $0 \leq \Re(s) \leq 1$, lie on the critical line for $\pi \neq \tilde{\pi}$.

 Combining CASE 1 and CASE 2, we conclude that Theorem 7 holds as a specific case of Theorem 4 with $\lambda=\pi, \bar\lambda = \tilde{\pi}, k=1$.
\end{proof}

Actually, from the above proofs of Theorem 5, Theorem 6, and Theorem 7, we can note that each proof does not depend on the specific definition of the $L$-function $L(\lambda, s)$, but rather relies on the following general properties of $\Lambda(\lambda, s)$ and $L(\lambda, s)$: \\
\textbf{P1:} Symmetric functional equation between $\Lambda(\lambda, s)$ and $\Lambda(\bar{\lambda}, k-s)$: $\Lambda(\lambda, s) = \varepsilon(\lambda) \Lambda(\bar{\lambda}, k-s)$;\\
\textbf{P2:} $\Lambda(\lambda, s)$ or $s^r(k-s)^r\Lambda(\lambda, s)$ ($r\geq 1$) is an entire function of order 1, with Hadamard product expression;\\
\textbf{P3:} The complex zeros in the relevant Hadamard products appear in pairs $(\rho_i, \bar\rho_i)$ with multiplicity $m_i\ge 1$;\\
\textbf{P4:} All zeros in the relevant Hadamard products lie in the critical strip $0 \leq \Re(\rho) \leq k$ (guaranteed by the Euler product property) and satisfy $\displaystyle\sum_{\rho\neq 0, k} \frac{1}{|\rho|^2} < \infty$;\\
\textbf{P5:} The disjointness of real and complex non-trivial zero sets;\\
\textbf{P6:} The zeros of $\Lambda(\lambda, s)$ are precisely the non-trivial zeros of $L(\lambda, s)$.\\

 \textbf{P1} and \textbf{P3} guarantee that $m_i$ is the multiplicity of the $i$-th quadruple of zeros ($\rho_i, \bar\rho_i, k-\rho_i, k-\bar\rho_i$).\\

Therefore, we have the following result on the non-trivial zero distribution of all kinds of $L$-functions.\\

\textbf{Theorem 8:} The non-trivial zeros of any $L$-function in the critical strip $0 \leq \Re(s) \leq k$ lie on the critical line if only the properties \textbf{P1}, \textbf{P2}, \textbf{P3}, \textbf{P4}, \textbf{P5}, and \textbf{P6} are satisfied.

\begin{proof}
It is not difficult to see that \textbf{P1}, \textbf{P2}, \textbf{P3}, \textbf{P4}, and \textbf{P5} cover all the conditions in Theorem 4. Thus, we know by Theorem 4 that all the zeros, both real (if any) and  complex, of $\Lambda(\lambda, s)$ in the critical strip $0 \leq \Re(s) \leq k$ lie on the critical line. Furthermore, according to \textbf{P6}, we conclude that all the non-trivial zeros, both real (if any) and  complex, of $L(\lambda, s)$ in the critical strip $0 \leq \Re(s) \leq k$ lie on the critical line.\\
That completes the proof of Theorem 8.
\end{proof}

\textbf{Remark:} Conditions \textbf{P1-P3} and \textbf{P6} are well-known general properties of $L$-functions, see Chapter 5 of Ref.[5]. Condition \textbf{P4} follows from Theorem 9, which is an extended result of Lemma 5.5. In detail, $0 \leq \Re(\rho) \leq k$ reduces to the Euler product, while $\displaystyle\sum_{\rho\neq 0, k} \frac{1}{|\rho|^2} < \infty$ is a general result of complex analysis for entire functions of order 1. Condition \textbf{P5} is automatically satisfied by our definitions of $\mathcal{Z}_{r}$ (the set of zeros with $\Im(\rho)=0$) and $\mathcal{Z}_{c}$ (the set of zeros with $\Im(\rho)\neq 0$).
Therefore, Theorem 8 with conditions \textbf{P1-P6} is expected to serve as a guideline for constructing new zeta- or $L$-type functions.\\

\textbf{Theorem 9:} Let $L(\lambda, s)$ be an $L$-function,  \( \Lambda(\lambda, s) \), the corresponding completed \( L \)-function satisfying a functional equation of the form \( \Lambda(\lambda, s) = \epsilon(\lambda) \Lambda(\bar{\lambda}, k-s)\), with \( k\in \mathbb{R}, k > 0 \). Then all zeros \( \rho \) of \( \Lambda(\lambda, s) \) lie in the critical strip \( 0 \leq \Re(\rho) \leq k \). Moreover, for any \( \epsilon > 0 \), we have
\[
\sum_{\rho \neq 0, k} |\rho|^{-1-\epsilon} < \infty.
\]
where $\lambda$ denotes a mathematical object, $\overline\lambda$ is the dual of $\lambda$; $\varepsilon(\lambda)$ is a complex number of absolute value 1, called the "root number" of $L$-function $L(\lambda, s)$.
\begin{proof}
Defining \( s = kt \) and \( F(\lambda, t) = \Lambda(\lambda, k t) \), the functional equation transforms as follows:
\[
F(\lambda, t) = \Lambda(\lambda, k t) = \epsilon(\lambda) \Lambda(\bar{\lambda}, k - k t) = \epsilon(\lambda) \Lambda(\bar{\lambda}, k(1-t)) = \epsilon(\lambda) F(\bar\lambda, 1-t).
\]
We thus obtain a new entire function \( F(\lambda, t) \) with the standard functional equation \( F(\lambda, t) = \epsilon(\lambda) F(\bar{\lambda}, 1-t) \), to which Lemma 5.5 applies.

Let \( \tau \) be a zero of \( F(\lambda, t) \). By definition, \( \rho = k \tau \) is then a zero of \( \Lambda(\lambda, s) \), and this correspondence is bijective. By Lemma 5.5 we know that \( 0 \leq \Re(\tau) \leq 1 \), \( \sum_{\tau\neq 0, 1} |\tau|^{-1-\epsilon} < \infty \) (for any \( \epsilon > 0 \)), which immediately implies \( 0 \leq \Re(\rho) \leq k \), \( \sum_{\rho\neq 0, k} |\rho|^{-1-\epsilon} < \infty \) (for any \( \epsilon > 0 \)).

That completes the proof of Theorem 9.
\end{proof}


\section*{Acknowledgments}

I wish to express my deepest gratitude to my late supervisor, Professor Zhifang Zhang (January 15, 1931 – July 6, 2026). His teachings and kindness will always remain with me.

Special gratitude is extended to the PREPRINT platform preprints.org for making this work permanently available and citable.

\textbf{We should stand on the shoulders of our predecessors, not grovel at their feet.}

\textbf{Data Availability Statement:} Data availability is not applicable to this article as no data were created or analyzed in this research.

\textbf{Funding Statement:} No funding received for this research.

\textbf{Conflict of Interest statement:} No conflict of interest.
\bibliographystyle{plainnat}
\bibliography{Mynatbib} % Your .bib file

[1] Kenneth Hoffman, Ray Kunze (1971), Linear Algebra, Second Edition,  Prentice-Hall, Inc., Englewood Cliffs, New Jersey.

[2] Karatsuba A. A., Nathanson M. B. (1993), Basic Analytic Number Theory, Springer, Berlin, Heidelberg.

[3] Conway, J. B. (1978), Functions of One Complex Variable I, Second Edition,  New York: Springer-Verlag.

[4] Olaf Helmer (1940), Divisibility properties of integral functions, Duke Mathematical Journal, 6(2): 345-356.

[5] Iwaniec, H., Kowalski, E. (2004). Analytic Number Theory, American Mathematical Society Colloquium Publications, Vol. 53.

[6] Cohen, Henri (2007), Number theory, Volume II: Analytic and modern tools, Graduate Texts in Mathematics, 240, New York: Springer, doi:10.1007/978-0-387-49894-2, ISBN 978-0-387-49893-5.

[7] Cogdell, J. W. (2014). L-functions and functoriality. In Jianya Liu (Ed.), Automorphic forms and L-functions. Beijing: Higher Education Press.

[8] Iwaniec H., Sarnak  P. (2000), Perspectives on the analytic theory of L-functions. GAFA 2000 (Tel Aviv, 1999), Geom. Funct. Anal. 2000, Special Volume, Part II, 705-741

[9] Yujiao Jiang, Guangshi Lv (2023), A survey on coefficients of automorphic L-functions, Advances in Mathematics (China), 52(2): 201-223.

[10] H. Jacquet, J. Shalika (1976), A nonvanishing theorem for zeta-functions of $GL_n$, Invent. Math., 38, 1-16.

[11] S. Gelbart, F. Shahidi (2001), Boundedness of automorphic L-functions in vertical strips, J. Amer. Math. Soc. 14 (1), 79–107

[12] Peter Sarnak (2004), Nonvanishing of L-Functions on $\Re(s) = 1$, in Contributions to Automorphic Forms, Geometry Number Theory, editors Haruzo Hida, Dinakar Ramakrishnan, and Freydoon Shahidi, The John Hopkins University Press, Baltimore, 2004, 719–732.

\end{document}

V14 SUBMISSION - DETAILS OF CHANGES: 1. The title of the paper is modified to be more specific; 2. The subsection Modular form L-function is removed; 3. Some refinements in the proofs of Theorem 3-Theorem 8. 2026-03-31晚上
V13 SUBMISSION - DETAILS OF CHANGES: 1. The Abstract is modified; 2. The Section "Introduction" is removed; 3. The Section "Preliminary Lemmas" is added; 4: The Section "Appendix: Lemmas" is removed; 5. Some typos are corrected, some Remarks are added. 2026-03-18

In addition, Theorem 4 could be expected to serve as a blueprint for engineering new zeta- or $L$-type series (functions): once a symmetric functional equation and Hadamard product (featuring complex conjugate zeros pairing with uniquely determined multiplicities) are established through analytic continuation and completion (to get completed $L$- functions as entire functions of order 1), all non-trivial zeros are forced onto the critical line $\Re(s)=\frac{k}{2}$ ($k> 0$ is a real constant). For more details, see Section 2 and Section 3.

The sole purpose of \textbf{P6} is to link the zeros of $\Lambda(\lambda, s)$ with the non-trivial zeros of $L(\lambda, s)$; further details about \textbf{P6} can be found in Ref.[5] (p.96), with $f$ therein replaced by $\lambda$. 